\documentclass[11pt]{article}
\usepackage[margin=1in]{geometry}
\usepackage{amsmath,amssymb,booktabs,hyperref}
\title{GapTheo: A Reproducible Visual Atlas of Three-Gap Dynamics}
\author{Felipe Santibáñez-Leal}
\date{September 2026}

\begin{document}
\maketitle

\begin{abstract}
GapTheo is a browser-native research instrument for the three-gap theorem.
It keeps the direct finite rotation orbit as an authoritative oracle while
synchronizing exact gap certificates with Farey geometry, continued fractions,
dual return times, cyclic words, lattice coordinates, and a finite
zero-dimensional Rips filtration. The contribution is an expository and
reproducibility protocol, not a new theorem or proof.
\end{abstract}

\section{Scope and conventions}
For an irrational $\alpha\in(0,1)$, consider the points
\[
  \{0\alpha\},\{1\alpha\},\ldots,\{(N-1)\alpha\}\in\mathbb{R}/\mathbb{Z}.
\]
After sorting the points, the consecutive circular arcs have at most three
distinct lengths. If three lengths occur, the largest equals the sum of the
other two. GapTheo calls $N$ the number of points; sources using $n+1$ points
therefore use $n=N-1$ in the corresponding formula.

The app separates four statuses: certified, boundary, contrast, and warning.
Exact rational inputs are periodic boundary cases. A finite decimal cannot
certify irrationality. Farthest-point and random allocators are contrast
processes, not theorem inputs.

\section{Direct orbit and gap genealogy}
The live lane computes the fractional parts directly, sorts them, and emits
each circular gap with endpoints, length, rank, and stable identifier. The
incremental view inserts points in time order. A new point lies in one parent
empty arc and emits two child arcs. This lineage makes the additive relation
visible while the sorted oracle remains the source of truth.

\section{Farey event atlas}
Following Hamada, the $\alpha$--height diagram is partitioned by the lines
$h=\{i\alpha\}$. In the relevant Farey cell
$a/b<\alpha<c/d$, the candidates are
\[
 b\alpha-a,\qquad c-d\alpha,\qquad
 (b-d)\alpha+c-a,
\]
with multiplicities $N-b$, $N-d$, and $b+d-N$ when the last is positive.
The third expression is the sum of the first two. GapTheo renders the active
cell, the moving $\alpha$ cursor, and residual comparisons with the direct
oracle.

\section{Dual returns}
For a target interval $[0,\beta)$, GapTheo selects indices satisfying
$\{i\alpha\}<\beta$ and records successive index differences. This dual
view has its own two-gap and three-gap regimes and an explicit strict-boundary
control. It is shown as related structure, not silently merged with the
primal theorem.

\section{Words, lattices, exchanges, and topology}
The cyclic sequence of gap labels becomes a word over at most three letters.
The app supplies a local symmetry probe around a largest-gap anchor. The
lattice view follows the space-of-lattices perspective: active short vectors
are shown alongside the circle certificate. Rotation is also rendered as a
two-interval exchange, with a zippered-rectangle-inspired schematic.

For the topology lens, a threshold grows on the sampled circle. In the finite
zero-dimensional Rips filtration, adjacent gaps are the merge thresholds.
GapTheo draws the component count and finite bars with direct gap identifiers.
It does not claim a general persistent-homology engine.

\section{Controlled extensions}
The farthest-point allocator places each new point at the midpoint of a
largest empty arc. Seeded random placement tests reproducibility. A
two-frequency sequence tests the change in the number of observed lengths.
These modes are useful falsification and boundary experiments. They do not
extend the one-dimensional theorem.

\section{Implementation and reproducibility}
The Python reference lane emits JSON certificates with canonical hashes. The
TypeScript live lane uses the same scenario and certificate vocabulary. The
analytic pipeline has no training stage: the template slot is explicitly not
applicable. The twelve checked-in scenarios cover quadratic and
transcendental-like irrational declarations, phase and finite-size changes,
near-rational conditioning, two exact rational periods, and two deterministic
allocator contrasts. The benchmark compares the three-gap bound, partition
closure, additive residual, rational collision labelling, dual return bound,
star discrepancy, and seeded reproducibility. A deployment validates this
committed evidence but never regenerates it.

\section{Discrepancy and finite-size event maps}
For sorted points $0\leq x_1\leq\cdots\leq x_N<1$, GapTheo reports the finite
star discrepancy
\[
 D_N^*=\max_i\left\{\left|\frac{i}{N}-x_i\right|,
 \left|x_i-\frac{i-1}{N}\right|\right\}.
\]
It also recomputes the number of distinct gaps for every $3\leq n\leq N$.
This event map makes continued-fraction scales visible without treating the
finite plot as an asymptotic theorem. Both views are derived from the same
direct partition oracle as the circle display.

\section{Novelty boundary}
The reviewed literature contains the individual mathematical mechanisms:
Farey and continued-fraction proofs, dual return times, combinatorics on words,
space of lattices, interval exchanges, and topology applications. GapTheo's
original contribution is their synchronized, deterministic, browser-native
visual atlas with an explicit proof boundary and reproducible artifacts. No
new theorem, proof, asymptotic claim, or formal verification result is
asserted.

\section{References}
\begin{enumerate}
\item T. Hamada, \emph{A concise geometric proof of the three distance theorem},
arXiv:2308.11999, \url{https://arxiv.org/abs/2308.11999}.
\item V. Alessandri and V. Berthé, \emph{Three distance theorems and
combinatorics on words}, \url{https://www.irif.fr/~berthe/Articles/3d.pdf}.
\item V. Berthé and C. Reutenauer, \emph{On the Three-Distance Theorem},
HAL:hal-04769002v1, \url{https://hal.science/hal-04769002v1}.
\item J. Marklof and A. Strömbergsson, \emph{The three gap theorem and the
space of lattices}, \url{https://arxiv.org/abs/1612.04906}.
\item A. Taha, \emph{The Three Gap Theorem, Interval Exchange Transformations,
and Zippered Rectangles}, \url{https://arxiv.org/abs/1708.04380}.
\item A. Dasgupta and D. Roeder, \emph{Symmetries of the Three-Gap Theorem},
\url{https://doi.org/10.1080/00029890.2022.2158021}.
\item L. Suarez Salas and J. A. Perea, \emph{Estimation of Persistence
Diagrams via the Three Gap Theorem}, \url{https://arxiv.org/abs/2603.04570}.
\end{enumerate}

\end{document}
